\documentclass[10pt,a4paper]{amsart}
\usepackage[margin=19mm]{geometry}
\usepackage{amsmath,amssymb,mathtools}
\usepackage[scr=rsfs,cal=euler]{mathalfa}
\usepackage{xfrac}
\usepackage[unicode,hidelinks,bookmarks=false]{hyperref}
\newcommand{\ie}{i.e.\ }
\newcommand{\cf}{cf.\ }
\newcommand{\resp}{respectively\ }
\newcommand{\Subsection}{\subsection}
\providecommand{\nobreakdash}{\nobreak\hbox{-}}
\setlength{\emergencystretch}{2em}
\multlinegap0pt
\usepackage{mathtools}
\usepackage{amssymb}
\usepackage[shortlabels]{enumitem}
\newtheorem{lemma}{Lemma}
\usepackage{cleveref}
\crefname{lemma}{Lemma}{Lemmas}
\title{Exact extremal constructions for the inducibility of blowup graphs}
\author{Wanfang Chen, Xizhi Liu}
\date{Source: arXiv:2606.06202v1 (CC BY 4.0)}
\begin{document}
\maketitle

\noindent\emph{Source and licence.} Exact extremal constructions for the inducibility of blowup graphs. Version arXiv:2606.06202v1 (CC BY 4.0), distributed by arXiv under the Creative Commons Attribution 4.0 International licence, http://creativecommons.org/licenses/by/4.0/, as recorded in the arXiv licence metadata for this exact version and shown on the abstract page https://arxiv.org/abs/2606.06202v1. Full licence notice: \texttt{licenses/arxiv-2606.06202-CC-BY-4.0.txt}.\par\smallskip

\noindent\textit{Editorial scope.} Counted unit: the article's lemma lem:many-forced-errors (source0.tex lines 602--609). The article's complete original proof of it is delivered with it (source0.tex lines 611--652, one \texttt{proof} environment of 42 lines). No mathematics is written, repaired or re-derived: the statement and the proof are the article's own lines, in the article's own notation, and no retained line is substituted or re-flowed.  Retained uncounted as the source-local context the proof needs: lines 210--213.  Nothing was omitted from any retained span, so the package carries no omission record.  No retained line contains a citation, so the wrapper carries no bibliography and no bibliography entry.  No retained line contains a figure environment, an image-inclusion command or drawing code, so no image is delivered and the package declares no asset dependency.  Editorial formatting only, in addition: an AMS \texttt{amsart} preamble that restates the article's own declarations and macros used by the retained lines, namely the environments \texttt{lemma} on separate counters, together with the standard packages those lines need.  The retained environments print in this document as Lemma 1 in order of appearance, whereas the article prints the same items with the numbering of its own full document.  The complete native source of the article as distributed in the arXiv e-print for this exact version is bundled unchanged as \texttt{sources/202169/source0.tex}.
\begin{lemma}
\label{lem:twin-free-auto}
If $R$ is twin-free, then every strong homomorphism $R\to R$ is an automorphism.
\end{lemma}

\begin{lemma}
\label{lem:many-forced-errors}
Let $H=R(\mathbf k)$, where $R$ is twin-free and $|V(R)|\ge2$. There is a constant $c_*=c_*(H)>0$ such that the following holds for every sufficiently large $h$. Put $F\coloneqq H^{(h)}=R(h\mathbf k)$. If $\theta:V(F)\to V(R)$ is not a strong homomorphism, then there are $a_0\in V(F)$ and $U\subseteq V(F)\setminus\{a_0\}$ with $|U|\ge c_*h$ such that
\[
        A_F(a_0,a)\ne A_R(\theta(a_0),\theta(a))
        \qquad\text{for all }a\in U.
\]
\end{lemma}

\begin{proof}
Let $r\coloneqq |V(R)|$, and let $k_{\min}\coloneqq\min_x\mathbf k(x)$. For each $x\in V(R)$, choose a value $\sigma(x)\in V(R)$ which occurs most often among the values of $\theta$ on the type class $\pi^{-1}(x)$, and set
\[
        M_x\coloneqq\{a\in\pi^{-1}(x):\theta(a)=\sigma(x)\}.
\]
Then
\[
        |M_x|\ge \frac{h\mathbf k(x)}{r}\ge \frac{hk_{\min}}{r}
        \qquad\text{for every }x\in V(R).
\]

First suppose that $\sigma:R\to R$ is not a strong homomorphism. Choose $x,y\in V(R)$ such that
\[
        A_R(x,y)\ne A_R(\sigma(x),\sigma(y)).
\]
Here $x\ne y$, since both graphs are loopless. Pick any $a_0\in M_x$ and take $U=M_y$. For every $a\in U$,
\[
        A_F(a_0,a)=A_R(x,y)
        \ne A_R(\sigma(x),\sigma(y))
        =A_R(\theta(a_0),\theta(a)),
\]
and $|U|\ge hk_{\min}/r$.

It remains to consider the case where $\sigma$ is a strong homomorphism. Since $R$ is twin-free, \cref{lem:twin-free-auto} implies that $\sigma$ is an automorphism. If $\theta(a)=\sigma(\pi(a))$ for every $a\in V(F)$, then $\theta$ would be a strong homomorphism, contrary to the assumption. Hence there is $a_0\in V(F)$ such that $\theta(a_0)\ne\sigma(\pi(a_0))$. By twin-freeness, there is $y\in V(R)$ such that
\[
        A_R(\theta(a_0),y)\ne A_R(\sigma(\pi(a_0)),y).
\]
Since $\sigma$ is onto, we may write $y=\sigma(w)$ for some $w\in V(R)$. Thus
\[
        A_R(\theta(a_0),\sigma(w))
        \ne A_R(\sigma(\pi(a_0)),\sigma(w))=A_R(\pi(a_0),w).
\]
Now take $U\coloneqq M_w\setminus\{a_0\}$. For every $a\in U$,
\[
        A_F(a_0,a)=A_R(\pi(a_0),w)
        \ne A_R(\theta(a_0),\sigma(w))
        =A_R(\theta(a_0),\theta(a)).
\]
Moreover, $|U|\ge hk_{\min}/r-1$.

Thus, in both cases, the desired conclusion holds with $c_*\coloneqq k_{\min}/(2r)$, after increasing the lower threshold on $h$ if necessary.
\end{proof}

\end{document}
